\section{Conclusion}
\label{sec:conclusion}

Two classical finite-injury priority arguments are formalized in Lean 4
over one explicitly finitary model of oracle computation.  The informal
temporal reasoning of the textbook proofs --- choose a fresh witness, wait
for convergence, restrain below the use, initialize what is weaker, argue
that the stronger requirements eventually go quiet --- becomes finite
executable stages, explicit state components, recorded uses, preserved
invariants, a bounded progress measure, and an induction along priority.
The constructions are not merely type-checked but executable: the stage
machines run, and produce the witnesses, restraints and enumerations the
proofs reason about.

The methodological lesson is about representation.  Priority arguments are
not beyond current provers, but they require representations that expose
finite information at exactly the points where the informal proof appeals
to it.  Recording the use inside the evaluator's result, rather than
recovering it after the fact, is what makes restraint a manipulable object;
keeping approximations as monotone finite sets rather than prefix-fixed
Boolean strings is what keeps the limit sets from being decidable and the
theorem from being false.

The second lesson is the one we did not expect, and it is negative.  A
foundation built for one priority argument served a second one completely
--- including a lemma whose importance in the second theorem exceeds its
importance in the first --- while the priority layer itself transferred not
at all, and no generalization of it appears to be worth stating.  ``Finite
injury'' names a family of arguments related by shape, not a lemma one
proves once.  A useful reusable priority framework, if there is one, will
have to abstract requirements, strategies, priority order, initialization
and progress measures at a level neither of these two developments
currently occupies; we regard that as open rather than as routine
engineering.

The natural next target is infinite injury.  It is what the omitted
conclusion of the Sacks Splitting Theorem needs --- lowness of the two
halves --- and it is the first place where the finite progress measure that
made the first theorem tractable is unavailable in principle rather than by
accident of design.  The other loose ends are smaller.  Transitivity of the
local reducibility relation, listed as future work in an earlier draft of
this development, is now proved by query substitution, so the relation is a
preorder; what remains is that the bridge to Mathlib's oracle model is
proved in one direction only, which suffices for the negative results here
but not for a general transfer, and that no degree structure --- the
quotient, the jump, completeness --- is formalized, so the classical
reading of these theorems as statements about degrees remains external to
the artifact.
